\documentclass[10pt,a4paper]{amsart}
\usepackage[margin=19mm]{geometry}
\usepackage{amsmath,amssymb,mathtools}
\usepackage[scr=rsfs,cal=euler]{mathalfa}
\usepackage{xfrac}
\usepackage[unicode,hidelinks,bookmarks=false]{hyperref}
\newcommand{\ie}{i.e.\ }
\newcommand{\cf}{cf.\ }
\newcommand{\resp}{respectively\ }
\newcommand{\Subsection}{\subsection}
\providecommand{\nobreakdash}{\nobreak\hbox{-}}
\setlength{\emergencystretch}{2em}
\multlinegap0pt
\usepackage{bm}
\usepackage{bbm}
\usepackage{dsfont}
\usepackage{xspace}
\usepackage{cancel}
\usepackage{mathtools}
\usepackage{amsthm}
\makeatletter
\@ifundefined{algorithm}{\newtheorem{algorithm}{Finite element Algorithm}}{}
\@ifundefined{theorem}{\newtheorem{theorem}{Theorem}}{}
\@ifundefined{lemma}{\newtheorem{lemma}[theorem]{Lemma}}{}
\makeatother
\def\E{{\mathcal{E}}}
\def\T{{\mathcal{T}}}
\def\b0{\boldsymbol{0}}
\def\bf{{\mathbf{f}}}
\def\bh{{\mathbf{H}}}
\def\bi{{\mathbf{I}}}
\def\bn{{\mathbf{n}}}
\def\bu{{\mathbf{u}}}
\def\bv{{\mathbf{v}}}
\def\bw{{\mathbf{w}}}
\def\sumT{\sum_{T\in\mathcal{T}_h}}     %new
\newcommand{\Real}{\mathbb{R}}
\newcommand{\di}{\text{div}}
\newcommand{\eps}{\varepsilon}
\newcommand{\la}{\langle}
\newcommand{\ra}{\rangle_{\partial T}}
\title[The Weak Galerkin and Crouzeix-Raviart element method for el]{The Weak Galerkin and Crouzeix-Raviart element method for elastic eigenvalue problems}
\author{Wei Lu, Hehu Xie,, Qilong Zhai}
\date{Source: arXiv:2508.02064v1 (CC BY 4.0)}
\begin{document}
\maketitle

\noindent\emph{Source and licence.} The Weak Galerkin and Crouzeix-Raviart element method for elastic eigenvalue problems. Version arXiv:2508.02064v1 (CC BY 4.0), distributed under the Creative Commons Attribution 4.0 International licence, as recorded in the arXiv licence metadata for this exact version and shown on the abstract page https://arxiv.org/abs/2508.02064v1. Full rights notice: licenses/arxiv-2508.02064-CC-BY-4.0.txt.\par\smallskip

\noindent\textit{Editorial scope.} Counted unit: the Theorem whose source label is \texttt{lb1} in the article (source0.tex lines 1209--1214, source label \texttt{lb1}), delivered together with the article's own complete original proof of it (source0.tex lines 1215--1272, one \texttt{proof} environment of 58 lines). No mathematics is written, repaired or re-derived: statement and proof are the article's own text line for line. Required context retained uncounted: vartion1 (lines 166--168); CR-int (lines 1125--1127); errint (lines 1130--1135); vartionh (lines 1154--1158); errH1 (lines 1187--1197). Every definition, display and label of those lines is retained unchanged. Omitted from this package and not redistributed: 1--165, 169--1124, 1128--1129, 1136--1153, 1159--1186, 1198--1208, 1273--1768. Nothing inside a retained excerpt is omitted or altered: every retained body is delivered exactly as the article's lines are written, so no omission receipt of any kind accompanies this package. Citations: the retained excerpts carry no citation command at all, so no bibliography entry of the article is transcribed and no reference is redistributed. Reference adaptation: none was needed and none was made; every reference inside the delivered unit resolves inside it, so no unresolved reference or citation is produced. Printed slips kept literal: no printed word, symbol, hypothesis or number of the counted statement or of its proof was altered. Layout and class: the article's own class is \texttt{siamltex704} with packages including amssymb, amsmath, graphicx, amscd, mathrsfs, color, xcolor, comment, dsfont, pifont, longtable, float, amsfonts, hyperref; the wrapper is the lane's pinned \texttt{amsart} editorial wrapper at 10pt on A4, which restates the theorem-like environments the retained text needs and adds the article's own macro definitions that the retained text needs (\texttt{\textbackslash E}, \texttt{\textbackslash Real}, \texttt{\textbackslash T}, \texttt{\textbackslash b}, \texttt{\textbackslash bf}, \texttt{\textbackslash bh}, \texttt{\textbackslash bi}, \texttt{\textbackslash bn}, \texttt{\textbackslash bu}, \texttt{\textbackslash bv}, \texttt{\textbackslash bw}, \texttt{\textbackslash di}, \texttt{\textbackslash eps}, \texttt{\textbackslash la}, \texttt{\textbackslash ra}, \texttt{\textbackslash sumT}), with extra wrapper packages (bm, bbm, dsfont, xspace, cancel, mathtools). The delivered numbering is the unit's own. Assets: no retained span contains a figure environment, a graphics-inclusion command, a PDF-page inclusion, a table or a TikZ picture; no image file is copied into this package, the package declares no asset dependency and no image file is redistributed with it. Licence: version arXiv:2508.02064v1 of this article is distributed under the Creative Commons Attribution 4.0 International licence, as recorded in the arXiv licence metadata for this exact version and shown on the abstract page https://arxiv.org/abs/2508.02064v1. A full rights notice is delivered at licenses/arxiv-2508.02064-CC-BY-4.0.txt. Characters: every retained excerpt is pure ASCII, so the delivered engine, XeLaTeX with its default Latin Modern text font, reports no missing character.
\begin{eqnarray}\label{vartion1}
	a(\bu,\bv)=\gamma b(\bu,\bv),\quad\forall\bv\in H_E^1(\Omega),
\end{eqnarray}

\begin{eqnarray}\label{CR-int}
	\int_{e}(\bv-\bi_h\bv)ds=\b0,\quad\forall e\in\E_h.
\end{eqnarray}

\begin{lemma}
	For any $\bv\in\bh^{1+s}(\Omega)$, there holds
\begin{align}\label{errint}
	\|\bv-\bi_h\bv\|+h|\bv-\bi_h\bv|_1\lesssim& h^{1+s}|\bv|_{1+s}.
\end{align}
\end{lemma}

\begin{algorithm}Find $\gamma_h\in\Real$ and $\bu_h\in V_h$ such that $b_h(\bu_h,\bu_h)$\\$=1$ and
	\begin{eqnarray}\label{vartionh}
		a_h(\bu_h,\bv_h)=\gamma_hb_h(\bu_h,\bv_h),\quad\forall\bv_h\in V_h,
	\end{eqnarray}
\end{algorithm}

\begin{lemma}
	Let $(\gamma,\bu)$ and $(\gamma_h,\bu_h)$ be the $j-$th eigenpair of \eqref{vartion1} and \eqref{vartionh}, respectively, then
	\begin{align}
		|\gamma-\gamma_h|\lesssim&\gamma(h)^{-2}h^{2s},
	\end{align}
	and there exists eigenfunction $\bu$ corresponding to $\gamma$ satisfies
	\begin{align}
		||\bu-\bu_h||_h\lesssim& \gamma(h)^{-1}h^s,\label{errH1}\\
		||\bu-\bu_h||\lesssim& \gamma(h)^{-2}h^{2s},\label{errL2}.
	\end{align}
\end{lemma}

\begin{theorem}
	Suppose that $\bu\in\bh^{1+s}(\Omega)$, we have
	\begin{align}\label{lb1}
		|a_h(\bu-\bi_h\bu,\bu_h)|\lesssim&h^{2s}.
	\end{align}
\end{theorem}

\begin{proof} It follows from the definition in \eqref{vartionh} that
	\begin{align*}
		a_h(\bu-\bi_h\bu,\bu_h)=&2\mu\sumT(\eps(\bu-\bi_h\bu),\eps(\bu_h))_T +\lambda\sumT(\di(\bu-\bi_h\bu),\di\bu_h)_T\\
		\quad&+\gamma(h)\sum\limits_{e\in\E_h^0}\frac{2\mu}{h_e}\langle[\![\bu-\bi_h\bu]\!],[\![\bu_h]\!]\rangle_e.
	\end{align*}
	By the Green formulation, we have
	\begin{align*}
		(\eps(\bu-\bi_h\bu),\eps(\bu_h))_T=&(\nabla(\bv-\bi_h\bv),\eps(\bu_h))_T
		\\
		=&\la\bv-\bi_h\bv,\eps(\bu_h)\bn\ra-(\bv-\bi_h\bv,\di(\eps(\bu_h)))_T,
	\end{align*}
	notice that $\eps(\bu_h)$ is a constant function, which together with \eqref{CR-int} verifies
	\begin{align*}
		(\eps(\bu-\bi_h\bu),\eps(\bu_h))_T=&0,\quad\forall T\in\T_h.
	\end{align*}
	Similarly, we get
	\begin{align*}
		(\di(\bu-\bi_h\bu),\di(\bu_h))_T=&0,\quad\forall T\in\T_h.
	\end{align*}
	Therefore, we have
	\begin{eqnarray}\label{ahM}
		a_h(\bu-\bi_h\bu,\bu_h)=\gamma(h)\sum\limits_{e\in\E_h^0}\frac{2\mu}{h_e}\langle[\![\bu-\bi_h\bu]\!],[\![\bu_h]\!]\rangle_e.
	\end{eqnarray}
	
	For $e\in\E_h^0$, define
	\begin{eqnarray*}
		P_e\bf=\frac{1}{|e|}\int_e\bf ds.
	\end{eqnarray*}
	
	 It is easy to check that $P_e[\![\bu_h]\!]=0$ and $[\![\bu]\!]=0$, then by Cauchy-Schwarz inequality, we derive
	\begin{align}\label{ahM2}
		|\langle[\![\bu-\bi_h\bu]\!],[\![\bu_h]\!]\rangle_e|=&
		|\langle[\![\bv]\!]-P_e[\![\bv]\!],[\![\bw]\!]\rangle_e|\nonumber
		\\
		=&|\langle[\![\bv]\!]-P_e[\![\bv]\!],[\![\bw]\!]-P_e[\![\bw]\!]\rangle_e|
		\\
		\lesssim&\big\|[\![\bv]\!]-P_e[\![\bv]\!]\big\|_e\big\|[\![\bw]\!]-P_e[\![\bw]\!]\big\|_e\nonumber,
	\end{align}
	where $\bv=\bu-\bi_h\bu$, $\bw=\bu-\bu_h$.
	
	Assume that $T^+,T^-\in\T_h$ such that $T^+\cap T^-=e$. Let $\bv^+=\bv|_{T^+},\bv^-=\bv|_{T^-}$. By the trace inequality and \eqref{errint}, we have
	\begin{align}\label{ahM3}
		\big\|[\![\bv]\!]-P_e[\![\bv]\!]\big\|_e=&
		\big\|(\bv^+-P_e\bv^+)-(\bv^--P_e\bv^-)\big\|_e\nonumber
		\\
		\lesssim&h^{\frac{1}{2}}|\bv|_{1,T^+\cup T^-}
		\\
		\lesssim& h^{\frac{1}{2}+s}|\bu|_{1+s,T^+\cup T^-}.\nonumber
	\end{align}
	
	Analogously, together with \eqref{errH1} we have
	\begin{align}\label{ahM4}
		\big\|[\![\bw]\!]-P_e[\![\bw]\!]\big\|_e
		\lesssim& \gamma(h)^{-1}h^{\frac{1}{2}+s}|\bu|_{1+s,T^+\cup T^-}.
	\end{align}
	
	Substituting \eqref{ahM2},\eqref{ahM3} and \eqref{ahM4} into \eqref{ahM} we get \eqref{lb1}, which completes the proof.
\end{proof}

\end{document}
